\section{Slopes, foliations, and polar parts}
\label{sec:slopes}

We retain the smooth projective threefold $V$ and the cycle $D$ of
Proposition~\ref{prop:prepared-cycle}.  Fix a very ample divisor $H$ on $V$.
The numerical properties used in this section are
\begin{equation}
  H^2D>0,\qquad HD^2=0,\qquad K_VDA=0
  \quad\text{for every divisor class }A.
  \label{eq:slope-numerical-input}
\end{equation}
We also use that $D$ is nef and that no curve of $D$-degree zero passes
through a very general point of $V$.  The ground field is uncountable,
algebraically closed, and of characteristic $p>3$.
We will obtain strong semistability and $c_2(V)D\geq0$, followed by
uniform bounds on vector-valued polar parts and the vanishing of scalar
polar parts.  The scalar vanishing starts the jet calculation in
Section~\ref{sec:jets}; the vector bounds control Laurent primitives in
Section~\ref{sec:laurent}.

\subsection{Slope calculus for the nef product}

For a torsion-free coherent sheaf $E$ of positive rank, set
\[
  \deg(E)=c_1(E)HD,\qquad
  \mu(E)=\frac{\deg(E)}{\rk(E)}.
\]
Stability and semistability will always refer to this slope.  A sheaf is
strongly semistable if all its iterated absolute Frobenius pullbacks are
semistable.  Although $D$ need not be ample, the elementary slope calculus
needed below still applies.  We give the details that are relevant here.

An effective divisor has nonnegative degree, since $H$ and $D$ are nef.
Saturating a subsheaf therefore does not decrease its degree.  Every
torsion-free sheaf $E$ embeds in a finite direct sum of copies of
$\cO_V(aH)$ for some $a$: dualize a sufficiently positive global-generation
map for $E^\vee$ and use $E\hookrightarrow E^{\vee\vee}$.  If $B\subset E$
has rank $b$, taking exterior powers and then determinants gives a nonzero
map from $\det B$ to some copy of $\cO_V(baH)$.  Hence
\[
  \mu(B)\le aH^2D.
\]
For ranks bounded by $\rk E$, the possible slopes form a discrete subset of
$\Q$.  Thus a maximal slope is attained.  Among saturated subsheaves of
maximal slope, take one of maximal rank.  The usual sum-and-intersection
argument shows that this subsheaf contains every other subsheaf of maximal
slope.  It is semistable, and the maximal slope of its torsion-free quotient
is strictly smaller.  Repeating on the quotient produces the
Harder--Narasimhan filtration.

Write $\mu_{\max}(E)$ and $\mu_{\min}(E)$ for the largest and smallest
slopes of its factors.  Additivity of degree and the image of a morphism
give the familiar implication
\begin{equation}
  \mu_{\min}(E)>\mu_{\max}(F)
  \quad\Longrightarrow\quad \Hom(E,F)=0.
  \label{eq:slope-hom-vanishing}
\end{equation}
Indeed a nonzero image is both a torsion-free quotient of the source and
a subsheaf of the target, which would force its slope to be at least the
left-hand extreme and at most the right-hand extreme.  Dualizing the
filtration on the open set where its terms and factors are locally free
gives
\[
  \mu_{\max}(E^\vee)=-\mu_{\min}(E),\qquad
  \mu_{\min}(E^\vee)=-\mu_{\max}(E).
\]
The omitted set has codimension at least two and does not change any
determinant or degree.  The same observation permits all ensuing bundle
calculations on such open sets, with saturation or reflexive extension
when returning to $V$.

Tensoring by a line bundle adds its degree to every slope.  More generally,
the same rule holds for a rank-one torsion-free factor after restriction
to the open set where it is invertible.  In particular,
Equation~\eqref{eq:slope-numerical-input} shows that twisting by $iD$ leaves
all slopes unchanged, for every $i\in\Z$.  We will use this rank-one rule;
no assertion about semistability of arbitrary tensor products is required.

\subsection{Connections and distributions of rank at most two}

Let $F_V\colon V\to V$ denote absolute Frobenius.  A Frobenius pullback
$F_V^*E$ has its canonical connection, which in a pulled-back local frame
differentiates the coefficients.  We first record the elementary descent
observation behind its use.

\begin{lemma}[The connection detects destabilization]
\label{lem:connection-destabilization}
Let $E$ be a semistable torsion-free sheaf.  If $F_V^*E$ is unstable and
$B$ is its first Harder--Narasimhan term, then the connection induces a
nonzero morphism on a big open set
\[
  B\longrightarrow (F_V^*E/B)\otimes\Omega_V^1.
\]
\end{lemma}

\begin{proof}
The induced map is $\cO_V$-linear by the Leibniz rule.  Suppose it vanishes.
Then the generic subspace defined by $B$ is horizontal.  Let $K=k(V)$ and
choose a pulled-back basis of the ambient generic vector space.  In a
Grassmannian chart, represent this subspace by a matrix $(I\mid M)$.
Horizontality says that every entry of $M$ has zero differential.  Since
$K$ is separably generated over the perfect field $k$, the kernel of
$d\colon K\to\Omega_{K/k}^1$ is $K^p$.  Taking $p$-th roots of the entries
of $M$ therefore defines a subspace before Frobenius pullback.

Saturate this subspace to a subsheaf $B_0\subset E$.  At every
codimension-one point the relevant torsion-free quotients are free over
the discrete valuation ring.  Frobenius is flat on the smooth variety,
so $F_V^*B_0$ is saturated there and agrees with $B$ there: they have the
same generic subspace.  Consequently
\[
  \deg(B)=p\deg(B_0),\qquad
  \mu(B_0)=\frac{\mu(B)}p>
  \frac{\mu(F_V^*E)}p=\mu(E),
\]
contradicting semistability of $E$.
\end{proof}

The general connection criterion is also stated in
\cite[Corollary~2.4]{langer2004}.  The argument above makes explicit why
it applies to the present nef slope.

\begin{lemma}[The first destabilizing term is a foliation]
\label{lem:destabilizing-foliation}
If $T_V$ is unstable, its first Harder--Narasimhan term $A$ is closed under
Lie brackets and under $p$-th powers of derivations.
\end{lemma}

This rank-one and rank-two argument follows the foliation criterion in
\cite[Lemma~2.9]{xuzhang2019}; we give the slope calculation for the
nef product used here.

\begin{proof}
Put $Q=T_V/A$ and $a=\mu(A)$.  By
Equation~\eqref{eq:slope-numerical-input}, $\mu(T_V)=0$, so
\[
  a>0,\qquad \mu_{\max}(Q)<a.
\]
The rank of $A$ is one or two.  Work where $A$ and $Q$ are bundles.  The
bracket modulo $A$ is an $\cO_V$-linear map
$\bigwedge^2A\to Q$.  It vanishes automatically in rank one.  In rank two,
$Q$ has rank one and slope $-2a$, whereas $\bigwedge^2A$ has slope $2a$;
Equation~\eqref{eq:slope-hom-vanishing} again makes the map zero.

Once bracket closure holds, the $p$-th power of derivations modulo $A$
defines an $\cO_V$-linear map
\begin{equation}
  F_V^*A\longrightarrow Q.
  \label{eq:p-curvature-obstruction}
\end{equation}
For completeness, the two restricted-Lie identities needed for this
assertion can be seen directly.  For local derivations $u,v\in A$, set
$w=tu+v$ with $t$ a formal scalar.  In characteristic $p$,
\[
  \frac{d}{dt}w^p
  =\sum_{j=0}^{p-1}w^juw^{p-1-j}
  =\operatorname{ad}(w)^{p-1}(u).
\]
The last equality follows from
$\binom{p-1}{j}\equiv(-1)^j\pmod p$.  Every coefficient of the right-hand
side lies in the bracket-closed module $A$.  The mixed coefficients of
$w^p-t^pu^p-v^p$, whose degrees in $t$ lie between $1$ and $p-1$, thus
belong to $A$.  This proves additivity modulo $A$.

For a local function $f$, expanding $(fu)^p$ as a differential operator
gives
\[
  (fu)^p-f^pu^p=\sum_{j=1}^{p-1}c_ju^j.
\]
The left side is a derivation.  At the generic point, if $u\ne0$, choose
$x$ with $u(x)\ne0$.  For a largest index $j\ge2$ with $c_j\ne0$, the
$j$-fold commutator with multiplication by $x$ is
$j!c_j(u(x))^j\ne0$, while it must vanish for a derivation.  Hence only
the $j=1$ term can remain.  The scalar discrepancy is therefore a
multiple of $u$, proving Frobenius linearity modulo $A$ and hence
Equation~\eqref{eq:p-curvature-obstruction}.

If $\rk A=1$, its source has slope $pa>a>\mu_{\max}(Q)$ and the map is
zero.  Suppose $\rk A=2$.  In this case
\[
  \mu(Q)=-2a,\qquad \mu_{\max}(\Omega_V^1)=2a.
\]
If $F_V^*A$ is semistable, its minimum slope is $pa>-2a$.  Otherwise let
$L$ be its first Harder--Narasimhan term and let $R=F_V^*A/L$; both have
rank one.  Lemma~\ref{lem:connection-destabilization} supplies a nonzero
map
\[
  L\longrightarrow R\otimes\Omega_V^1.
\]
Since $R$ has rank one, the slope rules give
\[
  \mu(L)-\mu(R)\le2a,\qquad
  \mu(L)+\mu(R)=2pa.
\]
It follows that
\[
  \mu_{\min}(F_V^*A)=\mu(R)\ge pa-a>-2a.
\]
Thus Equation~\eqref{eq:slope-hom-vanishing} makes
Equation~\eqref{eq:p-curvature-obstruction} zero in either case.
\end{proof}

\subsection{The inseparable quotient and its canonical divisor}

We spell out the quotient construction to keep track of the sign in the
canonical formula.  These constructions and their use for tangent
semistability also appear in \cite[Section~2]{xuzhang2019}.

\begin{lemma}[Quotient by the distribution]
\label{lem:foliation-quotient}
Let $A\subset T_V$ be a saturated bracket-closed and $p$-closed subsheaf
of rank $d$.  There is a finite purely inseparable map $\pi\colon V\to W$
to a normal projective variety, with a factorization
\[
  V\xrightarrow{\pi}W\xrightarrow{\rho}V^{(1)}
\]
of relative Frobenius.  On open sets whose complements have codimension
at least two, the canonical bundles satisfy
\begin{equation}
  \omega_V\simeq\pi^*\omega_W\otimes(\det A)^{\otimes(p-1)}.
  \label{eq:foliation-canonical}
\end{equation}
\end{lemma}

\begin{proof}
On an affine open with ring $R$, let
\[
  P=\{f\in R: u(f)=0\text{ for every local section }u\in A\}.
\]
The ring $P$ lies between $R^p$ and $R$.  Since $R$ is finite over $R^p$,
$P$ is finite over $R^p$ and $R$ is finite over $P$.  These invariant
rings glue to the asserted factorization of Frobenius; in particular both
maps are finite and their underlying maps of topological spaces are
homeomorphisms.  The ring $P$ is normal.  Indeed, an element of
$\Frac(P)$ integral over $P$ is integral over $R$, hence belongs to the
normal ring $R$, and all the derivations annihilate it.  Thus it belongs
to $P$.  The variety $W$ is projective because it is finite over the
projective variety $V^{(1)}$.

Remove the codimension-at-least-two loci where $A$ or $T_V/A$ fails to be
locally free.  Near any remaining point choose functions
$x_1,\ldots,x_d$ whose differentials restrict to a basis of $A^\vee$,
and choose the dual local derivations $u_1,\ldots,u_d$.  Thus
$u_i(x_j)=\delta_{ij}$.  Bracket closure implies $[u_i,u_j]\in A$; this
bracket annihilates every $x_l$ and is therefore zero.  Likewise
$u_i^p\in A$ annihilates all the $x_l$, so $u_i^p=0$.

These commuting nilpotent derivations give an explicit invariant-ring
presentation.  For one derivation $u$ with $u(x)=1$, the coefficient
$(p-1)!^{-1}u^{p-1}(f)$ lies in its kernel.  Subtract its product with
$x^{p-1}$ and repeat with successively smaller powers.  This writes every
$f$ uniquely as a polynomial of degree less than $p$ in $x$, with
coefficients in the kernel of $u$.  Applying this procedure successively
to the commuting $u_i$ gives
\begin{equation}
  R=\bigoplus_{0\le I_1,\ldots,I_d<p}P\,x^I
  \simeq
  P[X_1,\ldots,X_d]/(X_1^p-c_1,\ldots,X_d^p-c_d),
  \qquad c_i=x_i^p\in P.
  \label{eq:foliation-local-presentation}
\end{equation}
Uniqueness follows by applying the same successive highest derivatives
to a putative relation.  In particular, $\pi$ is locally free of degree
$p^d$ on this open set.

We may further remove the singular locus of $W$ and its inverse image;
normality and finiteness ensure that these removed sets also have
codimension at least two.  The presentation in
Equation~\eqref{eq:foliation-local-presentation} now yields the exact
sequence of vector bundles
\begin{equation}
  0\longrightarrow F_V^*(A^\vee)
  \longrightarrow\pi^*\Omega_W^1
  \longrightarrow\Omega_V^1
  \longrightarrow A^\vee\longrightarrow0.
  \label{eq:foliation-cotangent}
\end{equation}
Here the last map is evaluation on $A$.  To verify local freeness of the
first image, put $N=\langle dc_1,\ldots,dc_d\rangle$ inside
$R\otimes_P\Omega_{P/k}^1$.  The presentation gives
\[
  \Omega_{R/k}^1\simeq
  \bigl((R\otimes_P\Omega_{P/k}^1)/N\bigr)
  \oplus\bigoplus_{i=1}^dR\,dx_i,
\]
because the relations have no differential in the $X_i$ directions.
Smoothness of $R$ makes the first summand locally free of rank $3-d$.
Smoothness of $P$ then makes $N$ a locally free direct summand of rank
$d$ in the rank-three bundle $R\otimes_P\Omega_{P/k}^1$.  The displayed
$d$ generators of $N$ are a basis: locally their surjection from $R^d$
onto a free module of rank $d$ is an isomorphism.  This proves the
required exactness at every point of the chosen open set.

To identify its transition maps, let $x'_j$ be another such set of
functions and write it in the basis $x^I$ over $P$.  Taking $p$-th powers
and differentiating downstairs shows that, after pullback, the
coefficient of $dc_l$ in $d((x'_j)^p)$ is $u_l(x'_j)^p$.  On the other
hand, the coefficient of the $l$-th basis vector of $A^\vee$ in
$dx'_j|_A$ is $u_l(x'_j)$.  The first term of
Equation~\eqref{eq:foliation-cotangent} is therefore precisely
$F_V^*(A^\vee)$.  Taking determinants in that sequence gives
\[
  \pi^*\omega_W
  \simeq(\det A)^{-p}\otimes\omega_V\otimes\det A,
\]
which is Equation~\eqref{eq:foliation-canonical}.
\end{proof}

\subsection{Strong semistability of the tangent bundle}

\begin{proposition}
\label{prop:strong-semistability}
The tangent bundle $T_V$ is strongly semistable for the slope
$c_1(-)HD/\rk(-)$.  Moreover,
\begin{equation}
  c_2(V)D\ge0.
  \label{eq:c2-nonnegative}
\end{equation}
\end{proposition}

\begin{proof}
We first prove semistability by the foliation and rational-curve
argument of \cite[Lemma~2.10]{xuzhang2019}, with the quantitative
bend-and-break statement specified below.  Suppose that $A$ is the first
destabilizing term, with $a=\mu(A)>0$ and $d=\rk A$.  By
Lemmas~\ref{lem:destabilizing-foliation} and
\ref{lem:foliation-quotient}, it has a finite purely inseparable quotient
$\pi\colon V\to W$.  On the big smooth open set where
Equation~\eqref{eq:foliation-canonical} holds, put
$\Lambda=-\pi^*K_W$.  Its extension as a divisor class on $V$ is
\[
  \Lambda=(p-1)c_1(A)-K_V,
  \qquad
  \Lambda HD=(p-1)da>0.
\]
No global Cartier property of $K_W$ is needed: its degree will be taken
only on curves contained in this smooth open set.

Let $H^{(1)}$ and $D^{(1)}$ denote the corresponding divisor classes on
$V^{(1)}$, and set
\[
  H'=\rho^*H^{(1)},\qquad D'=\rho^*D^{(1)}.
\]
Then $H'$ is ample and Cartier, $D'$ is nef and Cartier, and
\[
  \pi^*H'=pH,\qquad \pi^*D'=pD.
\]
For a positive integer $t$, take sufficiently divisible positive
multiples of $H'$ and $H'+tD'$ which are very ample.  Their general
complete intersection is a smooth projective integral curve $C_t$ wholly
contained in the open set under consideration.  Here is the avoidance
point in this assertion.  The complement has dimension at most one.  A
general first divisor meets it in finitely many points and contains none
of its curve components; a general second divisor avoids those finitely
many points.  Bertini on the smooth open set, and irreducibility for the
successive very ample sections, give the required curve.

Projection formula, computed where $\pi$ is locally free, cancels both
the degree of $\pi$ and the chosen complete-intersection multiples from
the degree ratio.  It gives
\begin{align}
  \frac{D'C_t}{-K_WC_t}
  &=\frac{pD\,H(H+tD)}{\Lambda H(H+tD)} \notag\\
  &=\frac{pDH^2}{\Lambda H^2+t\Lambda HD}
    \longrightarrow0
    \qquad(t\longrightarrow\infty).
  \label{eq:bend-break-ratio}
\end{align}
The denominator is positive for all sufficiently large $t$.

We use the following quantitative form of bend-and-break
\cite[Theorem~1.1]{jlr2026}: if $Y$ is projective over an algebraically
closed field, $C\subset Y_{\mathrm{sm}}$ is an integral curve with
$K_YC<0$, and $N$ is a nef real Cartier divisor, then every closed point
of $C$ lies on a rational curve $\Gamma\subset Y$ satisfying
\[
  N\Gamma\le(\dim Y+1)\frac{NC}{-K_YC}.
\]
This statement is valid in arbitrary characteristic.  Choose $t$ so large
that the ratio in Equation~\eqref{eq:bend-break-ratio} is less than $1/4$.
Applied on $W$ with $N=D'$, the theorem gives, through every closed point
$x\in C_t$, a rational curve $\Gamma$ with
\[
  0\le D'\Gamma<1.
\]
As $D'$ is Cartier, this degree is an integer and hence is zero.

We choose $C_t$ to witness the very-general-point contradiction as well.
Fix the product of the two chosen very ample linear systems, and let
$\mathcal T$ be its nonempty open locus of smooth integral complete
intersections avoiding the bad locus above.  This parameter space is
irreducible.  Its universal curve $\mathcal C\to\mathcal T$ is smooth:
each fiber is a transverse complete intersection inside $W_{\mathrm{sm}}$,
so the relative Jacobian criterion applies.  Its evaluation map
$\mathcal C\to W$ is dominant.  Indeed the ambient incidence of points
and pairs of divisors through them is a product of projective bundles
over $W$, hence is irreducible and surjects onto $W$; restricting it to
the nonempty open parameter locus $\mathcal T$ leaves a dense open
subset, whose evaluation still has dense image.

Let $Z_1,Z_2,\ldots$ be proper closed subsets of $V$ outside whose union
no $D$-trivial curve passes.  The finite map $\pi$ has proper closed
images $\pi(Z_j)$.  For each $j$, the locus in $\mathcal T$ of curves not
contained in $\pi(Z_j)$ is the image of
$\mathcal C\setminus\operatorname{ev}^{-1}(\pi(Z_j))$ under the smooth
map $\mathcal C\to\mathcal T$.  It is open because smooth maps are open,
and nonempty because evaluation is dominant.  Over the uncountable field,
the countably many resulting nonempty opens in $\mathcal T$ have a common
$k$-point.  On its integral curve $C_t$, each $C_t\cap\pi(Z_j)$ is finite.
Choose a closed point $x$ outside their union.  Its unique inverse-image
point $y\in V$ lies outside every $Z_j$.

The reduced inverse image of the rational curve $\Gamma$ through $x$
contains a curve $\widetilde\Gamma$ through $y$, finite and surjective
over $\Gamma$.  Projection formula gives
\[
  pD\widetilde\Gamma
  =(\pi^*D')\widetilde\Gamma
  =\deg(\widetilde\Gamma/\Gamma)D'\Gamma=0.
\]
This contradicts the choice of $y$.  Thus $T_V$ is semistable.

Suppose next that $e\ge1$ is the first integer for which
$F_V^{e*}T_V$ is unstable.  Let $A'$ be its first Harder--Narasimhan term
and $Q'=F_V^{e*}T_V/A'$.  Applying
Lemma~\ref{lem:connection-destabilization} to the semistable sheaf
$F_V^{(e-1)*}T_V$ gives a nonzero map
\[
  A'\longrightarrow Q'\otimes\Omega_V^1,
  \qquad\text{equivalently}\qquad
  A'\otimes(Q')^\vee\longrightarrow\Omega_V^1
\]
on a big open set.  The cotangent bundle is semistable of slope zero,
because $T_V$ is.  Since $\rk A'+\rk Q'=3$, at least one tensor factor
in the second source has rank one.  If $\rk Q'=1$, semistability of
$A'$ and the rank-one twisting rule give the minimum slope of the source
as $\mu(A')-\mu(Q')$.  If $\rk A'=1$, the quotient $Q'$ need not be
semistable; use instead
\[
  \mu_{\min}((Q')^\vee)=-\mu_{\max}(Q')
\]
and twist by the rank-one factor $A'$.  In both cases,
\[
  \mu_{\min}\bigl(A'\otimes(Q')^\vee\bigr)
  =\mu(A')-\mu_{\max}(Q')>0.
\]
This contradicts Equation~\eqref{eq:slope-hom-vanishing}.  All Frobenius
pullbacks of $T_V$ are therefore semistable.

Finally apply Langer's Bogomolov inequality
\cite[Theorem~3.2]{langer2004}; see also the addendum
\cite{langeradd2004}.  This theorem is formulated for a nef tuple
$D_1,\ldots,D_{n-1}$ with numerically nonzero product, and for a strongly
semistable torsion-free sheaf gives
$\Delta(E)D_2\cdots D_{n-1}\ge0$, where
$\Delta(E)=2\rk(E)c_2(E)-(\rk(E)-1)c_1(E)^2$.
Here the tuple is $(H,D)$, whose product is nonzero since $H^2D>0$.
For the rank-three tangent bundle this yields
\[
  \bigl(6c_2(V)-2K_V^2\bigr)D\ge0.
\]
Equation~\eqref{eq:slope-numerical-input} gives $K_V^2D=0$ and proves
Equation~\eqref{eq:c2-nonnegative}.
\end{proof}

\subsection{Uniform section bounds}

\begin{lemma}
\label{lem:degree-zero-line-sections}
If a line bundle $L$ has $\deg L\le0$, then $h^0(V,L)\le1$.
\end{lemma}

\begin{proof}
A nonzero section has an effective zero divisor, so its degree is
nonnegative and hence must be zero.  Suppose that two independent
sections $s_0,s_1$ exist.  Remove the base locus of the resulting pencil.
Through every point $x$ of the remaining open set there is a member of
the pencil: take a nonzero linear combination vanishing at $x$.  Its
divisor has degree zero.  Each prime component $S$ through $x$ therefore
satisfies $DHS=0$, because all component degrees are nonnegative.

The restriction of $D$ to $S$ is nef and numerically trivial.  One way
to see the latter assertion, allowing $S$ to be singular, is to take a
projective resolution $\sigma\colon\widetilde S\to S$.  The class
$L_S=\sigma^*(D|_S)$ is nef and is orthogonal to the big and nef class
$A_S=\sigma^*(H|_S)$, with $A_S^2>0$.  The surface Hodge index theorem
\cite[Chapter~V, Section~1]{hartshorne1977}
makes the intersection form negative definite on $A_S^\perp$, whereas
nefness gives $L_S^2\ge0$.  Hence $L_S$ is numerically zero, and the
projection formula gives the assertion on $S$.  A projective integral
curve in $S$ through $x$, obtained for example as a component of a
sufficiently ample section through $x$, consequently has $D$-degree zero.
Choosing $x$ very general contradicts the prepared-cycle property.
\end{proof}

\begin{lemma}[A uniform Hom bound]
\label{lem:semistable-hom-bound}
If $G_1,G_2$ are torsion-free semistable sheaves of slope zero, then
\begin{equation}
  \dim_k\Hom(G_1,G_2)\le\rk(G_1)\rk(G_2).
  \label{eq:hom-rank-bound}
\end{equation}
\end{lemma}

\begin{proof}
First suppose both sheaves are stable.  For a nonzero map, its image has
slope at least zero as a quotient of $G_1$, and at most zero as a
subsheaf of $G_2$.  A positive-rank kernel would make the first inequality
strict by stability of $G_1$.  An image of rank smaller than $\rk G_2$
would make the second inequality strict by stability of $G_2$.  Thus
every nonzero map has full rank, and the two sheaves have the same rank,
say $n$.

The determinant of each map is a global section of
\[
  L=\det(G_2)\otimes\det(G_1)^{-1},\qquad\deg L=0.
\]
It is initially defined on the locally free open set and extends across
its codimension-at-least-two complement.  By
Lemma~\ref{lem:degree-zero-line-sections}, $H^0(V,L)$ has dimension at
most one.  If $f,g$ were linearly independent maps, fix a nonzero section
$\tau$ spanning this space and write
\[
  \det(sf+tg)=P(s,t)\tau
\]
for a homogeneous polynomial $P$ of degree $n$.  Since $f$ has full rank,
$P$ is nonzero; since $n>0$ and $k$ is algebraically closed, $P$ vanishes
at some point of $\PP^1(k)$.  The corresponding linear combination is a
nonzero map by independence, but has zero determinant, a contradiction.
Thus the Hom space between stable factors has dimension at most one.

Every slope-zero semistable sheaf admits a finite filtration by saturated
subsheaves with torsion-free stable factors of slope zero.  Indeed if it
is not stable, a smaller-rank equal-slope subsheaf can be saturated; its
degree cannot increase strictly because of semistability.  Both it and
the quotient are semistable of slope zero, and induction on rank finishes
the construction.  Apply the left-exact Hom sequences successively to
such filtrations of $G_1$ and $G_2$.  The dimension of the Hom space is at
most the sum of the dimensions between all pairs of stable factors.
Each summand is at most one, and the numbers of factors are at most the
respective ranks.  This proves Equation~\eqref{eq:hom-rank-bound}.
\end{proof}

\begin{proposition}[Polar-part bounds]
\label{prop:polar-bounds}
Let $G$ be a semistable vector bundle of slope zero.  For every integer
$i$,
\begin{equation}
  h^0\bigl(V,G(iD)\bigr)\le\rk G,
  \qquad
  h^0\bigl(V,G^{\otimes2}(iD)\bigr)\le(\rk G)^2.
  \label{eq:uniform-bundle-sections}
\end{equation}
For each fixed bundle $J=G$ or $J=G^{\otimes2}$, the increasing spaces
\begin{equation}
  H^0\bigl(V,J(iD)/J\bigr),\qquad i=1,2,\ldots,
  \label{eq:polar-part-system}
\end{equation}
have bounded dimension and eventually stabilize under their natural
injections.  In particular their elements have a uniform bound on their
actual pole orders along every prime component of $D$.  This pole bound
may depend on $J$; it is uniform in the cut $i$ and in the chosen section.
\end{proposition}

\begin{proof}
All the twists $G(iD)$ are semistable of slope zero.  Apply
Lemma~\ref{lem:semistable-hom-bound} first to $\cO_V,G(iD)$ and then to
$G^\vee,G(iD)$.  The identities
\[
  H^0\bigl(V,G(iD)\bigr)=\Hom\bigl(\cO_V,G(iD)\bigr),
  \qquad
  H^0\bigl(V,G^{\otimes2}(iD)\bigr)
    =\Hom\bigl(G^\vee,G(iD)\bigr)
\]
give Equation~\eqref{eq:uniform-bundle-sections}.  In particular the
second bound uses the two semistable bundles $G^\vee$ and $G(iD)$,
without needing semistability of $G^{\otimes2}$.

For $i>0$, the cohomology sequence of
$0\to J\to J(iD)\to J(iD)/J\to0$ gives
\[
  h^0\bigl(V,J(iD)/J\bigr)
  \le h^0\bigl(V,J(iD)\bigr)+h^1(V,J).
\]
The right side is bounded independently of $i$.  Effectivity of $D$
gives injections of the quotient sheaves as $i$ increases, hence
injections of their finite-dimensional spaces of global sections.
A bounded increasing sequence of integer dimensions is eventually
constant, so these spaces stabilize.

More explicitly, choose $i_0$ after stabilization.  Every global polar
part, at every later cut, is then the image of an element of
$H^0(V,J(i_0D)/J)$.  At the generic point of a prime component
$D_\alpha$ of multiplicity $m_\alpha$, a local bundle frame measures its
pole order coefficientwise.  Such an image has pole order at most
$i_0m_\alpha$ there.  This bound is independent of the later cut and of
the element chosen.
\end{proof}

\begin{corollary}[No scalar polar parts]
\label{cor:no-scalar-poles}
For every $i>0$,
\[
  H^0\bigl(V,\cO_V(iD)/\cO_V\bigr)=0.
\]
\end{corollary}

\begin{proof}
Apply Proposition~\ref{prop:polar-bounds} to $G=\cO_V$.  Suppose a
nonzero scalar polar part $\eta$ exists.  It has a genuine negative
valuation at the generic point of some component of $D$.  To justify
this detection, take local rational-function representatives.  If all
their component valuations were nonnegative, they would be regular by
normality of $V$; they already have no poles off $D$.  The quotient
section would then be zero.

For every $e\ge0$, raising the local representatives to their $p^e$-th
powers defines a global section
\[
  \eta^{p^e}\in
  H^0\bigl(V,\cO_V(p^eiD)/\cO_V\bigr).
\]
This operation is well defined on polar parts: two representatives
differing by a regular function have $p^e$-th powers differing by a
regular function.  At the selected component the negative valuation is
multiplied by $p^e$.  These pole orders are unbounded, contradicting the
uniform bound in Proposition~\ref{prop:polar-bounds}.
\end{proof}
